\input{preamble}

%% Content of slides %%

\begin{document}
    \begin{frame}[plain]
        \titlepage
    \end{frame}
    %
    \begin{frame}{Proceso Estocástico}
        \begin{definition}[Proceso Estocástico\cite{luisrincon2019}]
            Un proceso estocástico es una colección de variables aleatorias $\{X_t : t\in T\}$ parametrizada por un conjunto $T$, llamado espacio parametral, en donde las variables toman valores en un conjunto $S$ llamado espacio de estados.
        \end{definition}
    \end{frame}
    \begin{frame}{Cadena de Markov}
        \begin{definition}[Propiedad de Markov]
            Un proceso estocástico se dice que cumple la propiedad de Markov si para todo $t_1<t_2<...<t_n<t \in T$ y todo $x_1,x_2,...,x_n$ se cumple $P(X_t\leq x : X_{t_1}=x_1, X_{t_2}=x_2, ..., X_{t_n}=x_n) = P(X_t\leq x : X_{t_n}=x_n)$
        \end{definition}
    \end{frame}
    \begin{frame}{Movimiento Browniano}
        \begin{definition}[Movimiento Browniano\cite{luisrincon2019}]
        Un movimiento Browniano unidimensional de parámetro $\sigma^2$ es un proceso estocástico $\{B_t : t\geq 0\}$ con valores en $\mathbb{R}$ que cumple las siguientes propiedades:
        \begin{enumerate}
            \item $B_0=0$
            \item Las trayectorias son continuas.
            \item El proceso tiene incrementos independientes.
            \item Para cualesquiera tiempos $0\leq s<t$, la variable tiene incrementos $B_t-B_s$ independientes con distribución $\mathcal{N}(0,\sigma^2(t-s))$
        \end{enumerate}
        \end{definition}
    \end{frame}
    \begin{frame}{Bibliografía}
        \bibliography{BMBibTeX}
        \bibliographystyle{ieeetr}
    \end{frame}
\end{document}

    
